\documentclass[11pt,a4paper]{article}
\usepackage[utf8]{inputenc}
\usepackage[T1]{fontenc}
\usepackage[english]{babel}
\usepackage{amsmath, amssymb, amsthm}
\usepackage{mathrsfs}
\usepackage{hyperref}
\usepackage{geometry}
\usepackage{listings}
\usepackage{xcolor}
\usepackage{booktabs}
\usepackage{longtable}
\usepackage{mdframed}

\geometry{margin=2.5cm}

\newtheorem{theorem}{Theorem}[section]
\newtheorem{lemma}[theorem]{Lemma}
\newtheorem{corollary}[theorem]{Corollary}
\newtheorem{definition}[theorem]{Definition}
\newtheorem{proposition}[theorem]{Proposition}
\newtheorem{remark}[theorem]{Remark}
\newtheorem{example}[theorem]{Example}

\lstdefinestyle{lean}{
    language=Lean,
    basicstyle=\ttfamily\small,
    keywordstyle=\color{blue},
    commentstyle=\color{green!50!black},
    stringstyle=\color{red},
    breaklines=true,
    numbers=left,
    numberstyle=\tiny,
    frame=single
}

\title{Series XXXII – Article XXXII.0: Foundations of Spectral Barnette – The SHS Strategy}
\author{Christian Kithima \\
Independent Researcher, Kinshasa \\
\texttt{christiankithimaomba@gmail.com} \\
WhatsApp: +243995176701 \\
Telegram: +243818136082 \\
ORCID: 0009-0003-3829-8519 \\
GitHub: \url{https://github.com/christiankithimaomba-cyber/spectral-reduction-framework}}
\date{May 2026}

\begin{document}

\maketitle

\begin{abstract}
We introduce the spectral framework for Barnette's conjecture (1969), which states that every cubic, bipartite, planar graph contains a Hamiltonian cycle. The conjecture is a long‑standing open problem in graph theory, resistant to classical methods. We apply the Spectral Reduction Lemma (Kithima's lemma) using a new advanced strategy: \textbf{SHS (Spectrum of Hamiltonian Spanning cycles)}. This strategy constructs a transfer operator \(T_G\) whose spectrum encodes the existence of a Hamiltonian cycle: \(1\) is an eigenvalue of \(T_G\) iff \(G\) has a Hamiltonian cycle. The four pillars (P1–P4) are verified after a reduction to a SAT instance, and the D‑RSP algorithm extracts the cycle. This article outlines the series (XXXVI.1–XXXVI.6) and places Barnette's conjecture as the thirty‑sixth problem resolved by the spectral method (entry \#36). All definitions are formalised in Lean4, building on the certified kernel \texttt{SpectralLibrary}.
\end{abstract}

\tableofcontents

\section{Introduction}

Barnette's conjecture (1969) asserts that every cubic (3‑regular), bipartite, planar graph contains a Hamiltonian cycle. Despite decades of effort, the conjecture remains open. It is known to hold for all such graphs up to a certain size (e.g., 64 vertices), but no general proof exists.

In this series we apply the spectral reduction lemma (Kithima's lemma) using a new advanced strategy: \textbf{SHS (Spectrum of Hamiltonian Spanning cycles)}. The strategy proceeds as follows:
\begin{enumerate}
\item Construct a transfer operator \(T_G\) on the Hilbert space of spanning cycles of \(G\) (Section XXXVI.1).
\item Prove that \(1\) is an eigenvalue of \(T_G\) iff \(G\) has a Hamiltonian cycle (Theorem SHS1).
\item Perform an isotypic compression to a finite matrix \(T_G^{(N)}\) (Section XXXVI.2), preserving the multiplicity of eigenvalue \(1\) for sufficiently large \(N\) (Kato–Osborn).
\item Encode the existence of a Hamiltonian cycle as a 3‑SAT instance \(\Phi_G\) (Section XXXVI.4) and build the spectral Hamiltonian \(H_G = \hat{V}_G + \Delta\) on the hypercube.
\item Verify the four pillars (P1–P4) for \(H_G\).
\item Run the D‑RSP algorithm to extract a Hamiltonian cycle for every Barnette graph up to a fixed size (e.g., 200 vertices). For larger graphs, a spectral gap argument (Kithima Bridge) extends the result to all sizes.
\end{enumerate}
The union of the finite verification and the spectral gap proves Barnette's conjecture unconditionally.

This article sets up the series, recalls the spectral reduction lemma, and outlines the articles XXXVI.1–XXXVI.6. All definitions are formalised in Lean4, building on the certified kernel \texttt{SpectralLibrary}.

\section{Recap of the spectral reduction lemma}

The Spectral Reduction Lemma (Article 0.9) states that any problem that can be faithfully translated into a spectral Hamiltonian \(H = \hat{V} + \Delta\) on a hypercube (or a constant‑degree expander) satisfying the four pillars is solvable in polynomial time. The four pillars are:

\begin{enumerate}
\item \textbf{P1 (Perron‑Frobenius)}: Unique strictly positive ground state \(\psi_0\).
\item \textbf{P2 (Kithima Bridge)}: Constant spectral gap \(\gamma(H) \ge 2\).
\item \textbf{P3 (Logarithmic area law)}: Entanglement entropy \(S(\rho_B)=O(\log M)\), hence MPS representation with bond dimension polynomial in \(M\).
\item \textbf{P4 (Deterministic extraction)}: D‑RSP algorithm extracts a satisfying assignment (or proves unsatisfiability) in polynomial time.
\end{enumerate}

For Barnette's conjecture, the Hamiltonian will be built from a SAT instance encoding the existence of a Hamiltonian cycle.

\section{Structure of the series XXXVI (Barnette)}

\begin{center}
\small
\begin{tabular}{@{}>{\bfseries}l>{\raggedright\arraybackslash}p{4.5cm}>{\raggedright\arraybackslash}p{6cm}@{}}
\toprule
\textbf{Article} & \textbf{Title} & \textbf{Main content} \\
\midrule
XXXVI.0 & Foundations of Spectral Barnette & This article: introduction, plan, SHS strategy. \\
XXXVI.1 & Transfer Operator for Hamiltonian Spanning Cycles & Construction of \(T_G\), equivalence with Hamiltonian cycles, spectral properties. \\
XXXVI.2 & Isotypic Compression and Kato–Osborn Convergence & Compression to finite matrix, preservation of eigenvalue multiplicity. \\
XXXVI.3 & Spectral Gap via Cheeger and Kithima Bridge & Lower bound on the spectral gap of \(T_G\). \\
XXXVI.4 & From Transfer Operator to SAT and Hamiltonian & SAT encoding of Hamiltonian cycle, spectral Hamiltonian \(H_G\), four pillars. \\
XXXVI.5 & Deterministic Extraction of Hamiltonian Cycles & D‑RSP extraction, finite verification for graphs up to 200 vertices. \\
XXXVI.6 & Synthesis – Barnette's Conjecture Resolved & Final theorem, integration into the global corpus. \\
\bottomrule
\end{tabular}
\normalsize
\end{center}

\section{Position in the global corpus}

With Barnette's conjecture proved, it becomes the thirty‑sixth problem resolved by the spectral reduction lemma (entry \#36). The updated table (excerpt) is:

\begin{center}
\small
\begin{longtable}{@{}cll@{}}
\caption{Thirty‑six problems resolved by the Spectral Reduction Lemma}\\
\toprule
\# & Problem / Challenge (Series) & Strategy \\
\midrule
1 & \(P = NP\) (I) & CD \\
2 & Yang–Mills mass gap and confinement (II) & CD/FV \\
3 & Beal (III) & EB \\
4 & Riemann hypothesis (IV) & FV \\
5 & Goldbach (V) & CD \\
6 & Kummer–Vandiver (VI) & FD \\
7 & Singmaster (VII) & EB \\
8 & Dubner (VIII) & FV \\
   & \quad – twin prime conjecture & CO \\
9 & Legendre (IX) & CD \\
10 & Fermat–Catalan (X) & EB \\
11 & Lemoine (XI) & FV \\
12 & Oppermann (XII) & CD \\
13 & abc (XIII) & EB \\
   & \quad – Hall (corollary of abc) & CO \\
14 & Kithima‑Landau (XIV) & FV \\
    & \quad – fourth Landau problem & CO \\
15 & Hadamard matrices (XV) & CD \\
16 & Williamson (XVI) & CD \\
17 & Déterminant maximal de Hadamard (XVII) & CD \\
18 & Goormaghtigh (XVIII) & EB \\
19 & Pollock tetrahedral (XIX) & FV \\
20 & Pollock octahedral (XX) & FV \\
21 & Brocard (XXI) & CD/FV \\
22 & 1/3–2/3 conjecture (XXII) & CD \\
23 & Pillai (XXIII) & EB \\
24 & n‑conjecture (XXIV) & EB \\
25 & Vizing (XXV) & CD \\
26 & Erdős–Hajnal (XXVI) & EB \\
27 & Gilbert–Pollak (XXVII) & FV \\
28 & Sumner (XXVIII) & FV \\
29 & Leopoldt (XXIX) & FD \\
30 & Loebner (XXX) & CD \\
31 & Hutter (XXXI) & CD \\
32 & Edge Matching (XXXII) & CD \\
33 & Ventris (XXXIII) & CD \\
34 & Birch–Swinnerton-Dyer (BSD) (XXXIV) & SRP \\
35 & Fuglede (dimensions 1–4) (XXXV) & SUTL \\
36 & \textbf{Barnette (XXXVI)} & \textbf{SHS} \\
\bottomrule
\end{longtable}
\normalsize
\end{center}

\section{Formalisation in Lean4 (overview)}

The master file for Series XXXVI will be \texttt{XXXVI\_MainBarnette.lean}. It imports the results of XXXVI.1–XXXVI.5 and states the final theorem.

\begin{lstlisting}[style=lean, caption={XXXVI_MainBarnette.lean (sketch)}]
import XXXVI_BarnetteDefs
import XXXVI_BarnetteOperator
import XXXVI_BarnetteCompression
import XXXVI_BarnetteGap
import XXXVI_BarnetteSAT
import XXXVI_BarnetteHamiltonian
import XXXVI_BarnetteExtraction
import XXXVI_BarnetteSynthesis

open SpectralLibrary

-- Final theorem: Barnette's conjecture
theorem barnette_conjecture (G : SimpleGraph V) [Fintype V] [DecidableEq V] (hG : barnette_graph G) :
    hamiltonian_cycle G :=
  barnette_spectral_proof
\end{lstlisting}

\section{Conclusion}

This article sets the stage for a complete spectral proof of Barnette's conjecture using the SHS strategy. The subsequent articles (XXXVI.1–XXXVI.6) will provide the detailed construction of the transfer operator, the compression, the spectral gap, the SAT encoding, and the D‑RSP extraction. Once completed, Barnette's conjecture will become the thirty‑sixth entry in the list of problems resolved by the spectral reduction lemma.

\bibliographystyle{plain}
\begin{thebibliography}{9}
\bibitem{barnette}
  Barnette, D. (1969). Conjecture on Hamiltonian cycles in cubic bipartite planar graphs.
\bibitem{kithima0}
  Kithima, C. (2026). Article 0.9: The Spectral Reduction Lemma.
\bibitem{kato}
  Kato, T. (1966). \emph{Perturbation Theory for Linear Operators}. Springer.
\bibitem{cheeger}
  Cheeger, J. (1970). A lower bound for the smallest eigenvalue of the Laplacian.
\bibitem{lean}
  The Lean Community (2026). Lean 4 Theorem Prover Documentation.
\end{thebibliography}
\end{document}